\chapter{2013年「大学物理1」模拟 期中试题}
\date{2013年6月??日}{大学物理教学团队}{A0715011}
{未央物理讲师 \href{https://qm.qq.com/q/UPbGudx8cK}{\textbf{Axia}}}
{\qrcode{https://qm.qq.com/q/asJiytD5pC}}{「未央」HDU 物理营}
{\qrcode{http://weixin.qq.com/r/mhKYgFTEMQlOrRDj90eI}}{未央学社公众号}
\section{选择}
\begin{choice}[A]{}
某物体的运动规律为$\frac{\d v}{\d t}=-kv^2t$，式中的$k$为大于零的常数。当$t=0$时，初速为$v_0$，则速率$v$与时间$t$的函数关系是
\begin{tasks}(4)
    \task $\frac{1}{v}=\frac{1}{2}kt^2+\frac{1}{v_0}$
    \task $v=-\frac{1}{2}kt^2+v_0$
    \task $v=\frac{1}{2}kt^2+v_0$
    \task $\frac{1}{v}=-\frac{1}{2}kt^2+\frac{1}{v_0}$
\end{tasks}
\end{choice}
\begin{choice}[B]{}
如图所示，置于水平光滑桌面上质量分别为$m_1$和$m_2$ 的物体 $A$ 和$B$之间夹有一轻弹簧.首先用双手挤压$A$和$B$使弹簧处于压缩状态，然后撤掉外力，则在$A$和$B$被弹开的过程中
\begin{tasks}(2)
    \task 系统的动量守恒，机械能不守恒
    \task 系统的动量守恒，机械能守恒
    \task 系统的动量不守恒，机械能守恒
    \task 系统的动量与机械能都不守恒
\end{tasks}
\end{choice}
\begin{choice}[D]{}
    轮圈半径为$R$，其质量$M$均匀分布在轮缘上，长为$R$，质量为$m$的均质辐条固定在轮心和轮缘间，辐条共有$2N$根.今若将辐条数$N$减少$N$根但保持轮通过轮心，垂直于轮平面轴的转动惯量保持不变，则轮圈的质量为
\begin{tasks}(4)
    \task $\frac{N}{12}m+M$
    \task $\frac{N}{6}m+M$
    \task $\frac{2N}{3}m+M$
    \task $\frac{N}{3}m+M$
\end{tasks}
\end{choice}
\begin{choice}[D]{}
如图所示，一人造地球卫星到地球中心$O$对最大距离和最小距离分别是$R_A$和$R_B$.设卫星对应的角动量分别是$L_A$和$L_B$，动能分别是$E_{k_A}$和$E_{k_B}$，则应有
\begin{tasks}(4)
    \task $\begin{cases}L_B>L_A\\E_{k_A}>E_{k_B}\end{cases}$
    \task $\begin{cases}L_B>L_A\\E_{k_A}=E_{k_B}\end{cases}$
    \task $\begin{cases}L_B<L_A\\E_{k_A}>E_{k_B}\end{cases}$
    \task $\begin{cases}L_B=L_A\\E_{k_A}<E_{k_B}\end{cases}$
\end{tasks}
\end{choice}